% generated by roam.py handbook export -- не редактировать, править заметки
% 20261004192832-метод_брента.org (body 90ff8fd477b8)

\begin{method}[{Метод Брента}]\label[method]{note:1BCFB956-2053-49A8-ADDF-0D85D53F0610}
  Совмещает в себе \hyperref[note:4E950058-62DB-4EDA-A71D-98CACBBC41AA]{метод парабол} и \hyperref[note:96DED7D9-3918-45C5-9850-E05030941645]{метод золотого сечения}.

  Мы должны отслеживать значения оптимизируемой функции $f$ в 6 точках: $a,b,x,v,w,u$

  Тогда $[a;b]$~--- интервал, в котором ищем минимум на текущей итерации.
  Точки $x, w, v$~--- три лучшие найденные точки, $f(x) \le f(w) \le f(v)$ ($x$~--- текущее приближение минимума, $w$~--- вторая по значению, $v$~--- предыдущее $w$); по ним строится парабола, как в \hyperref[note:4E950058-62DB-4EDA-A71D-98CACBBC41AA]{методе парабол}.

  А точка $u$ будет либо вершина параболы по точкам $x,w,v$, либо точка золотого сечения наибольшего из отрезков $[a;x]$, $[x;b]$

  Парабола строится только в случае, если $x,w,v$~--- различны. И вершина параболы $u^{*}$ становится новой точкой $u$, только если:

  \begin{itemize}
    \item $u^{*} \in [a;b]$
    \item $u^{*}$ отстоит от $x$ не более чем на половину шага, сделанного до предыдущего (т.е. для итерации $t-2$).
  \end{itemize}

  Иначе точка $u$  определяется \hyperref[note:96DED7D9-3918-45C5-9850-E05030941645]{золотым сечением}.
\end{method}

\begin{remark}
  После вычисления $f(u)$ отрезок $[a;b]$ сужается, как в золотом сечении, а тройка $x, w, v$ обновляется так, чтобы $x$ осталась лучшей точкой. Вблизи $x^{*}$ работают параболы~--- сходимость \hyperref[note:9FA20DA5-3A95-4683-9CB5-C6C9A903C86A]{сверхлинейная}; золотое сечение гарантирует \hyperref[note:DF2405BF-8B31-47EB-8DEF-B1A88034717C]{линейную} в худшем случае.
\end{remark}
